% 111-03(111쪽) — 원 $x^2+y^2=1$ 에 내접한 직사각형 ABCD(A 왼쪽 위 · D 오른쪽 위 · B 왼쪽 아래 · C 오른쪽 아래)와 대각선 BD(동경 OD).
% 스캔 화질이 낮아 TikZ 로 다시 그림. 라벨은 원문에 있는 A·B·C·D·O·$x$·$y$·$1$·$-1$·$x^2+y^2=1$ 만 둔다.
% 원문에 없는 $\theta$ 표시나 점 C 의 좌표(답)는 그리지 않는다.
\begin{tikzpicture}[line width=0.5pt, x=1.15cm, y=1.15cm, >=stealth]
  \draw[->, line width=0.7pt] (-1.42,0) -- (1.52,0) node[below=1pt, font=\small] {$x$};
  \draw[->, line width=0.7pt] (0,-1.42) -- (0,1.5) node[left=1pt, font=\small] {$y$};
  \draw[line width=0.6pt] (0,0) circle (1);
  \draw[line width=0.55pt] (-0.84,0.545) -- (0.84,0.545) -- (0.84,-0.545) -- (-0.84,-0.545) -- cycle;
  \draw[line width=0.55pt] (-0.84,-0.545) -- (0.84,0.545);
  \node[above left=-3.5pt and -3pt, font=\small] at (-0.84,0.545) {$\pt{A}$};
  \node[below left=-3.5pt and -3pt, font=\small] at (-0.84,-0.545) {$\pt{B}$};
  \node[below right=-3.5pt and -3pt, font=\small] at (0.84,-0.545) {$\pt{C}$};
  \node[above right=-3.5pt and -3pt, font=\small] at (0.84,0.545) {$\pt{D}$};
  \node[below left=-1.5pt and -1pt, font=\small] at (0,0) {$\pt{O}$};
  \node[left=1pt, font=\small] at (0,1) {$1$};
  \node[below left=-3pt and -3pt, font=\small] at (0,-1) {$-1$};
  \node[above right=-3pt and -3pt, font=\small] at (1,0) {$1$};
  \node[below left=-1pt and 0pt, font=\small] at (-1,0) {$-1$};
  \node[right=1pt, font=\small] at (0.1,1.18) {$x^2+y^2=1$};
\end{tikzpicture}
